\section{Future outlook and roadmap}
\label{sec:outlook}

\subsection{Where the field stands}
Ten years after the first automated molten glass printers, the direct route has split into two mature but disjoint capabilities. Crucible-fed extrusion produces architectural-scale, transparent soda-lime objects at kilograms per hour with bead-scale resolution and has crossed from laboratory to a small studio-manufacturing business with a credible circular-materials story \citep{inamura2018,mit2024masonry,stern2024cgc}. Laser-filament deposition produces fused-silica optics and micro-components at millilitres per hour with nanometre surfaces and has crossed into a first commercial desktop product and the first real-time controlled process \citep{lgd2024homog,kth2024microlens,nobula_aniwaa,dgf2026control}. Between them lies unoccupied territory: no process prints transparent borosilicate or aluminosilicate at centimetre scale, no crucible-fed machine has real-time geometric control, no laser process exceeds a few cubic centimetres per hour, and no printed glass part of any kind has a published Weibull strength distribution or a quantitative three-dimensional residual-stress map. The roadmap below is organised around closing these gaps and is stated as a sequence of milestones that the evidence suggests are reachable within roughly five years.

\subsection{Near-term milestones (1--2 years)}
\begin{enumerate}[leftmargin=*,itemsep=2pt]
  \item \textbf{Geometric closed-loop control on crucible-fed printers.} Thermal imaging is already on G3DP3 \citep{stern2024cgc}; adding bead-width extraction from the thermal image and a slump observer based on the sag-velocity scaling of \cref{sec:physics} would convert the present between-print mass correction into within-layer travel-speed control. The 2026 digital-glass-forming controller shows the architecture (deterministic real-time platform, camera in the loop) \citep{dgf2026control}.
  \item \textbf{Standardised optical and mechanical reporting.} A minimal coupon set---polished-face spectral transmission, transmitted wavefront, scattered-light polariscopy, four-point bending with Weibull statistics in two orientations---printed alongside every published part would make the modalities comparable. The tools exist \citep{aben2008,inamura2018}; the practice does not.
  \item \textbf{Recycled-glass qualification.} The viscosity-matching workflow \citep{stern2024cgc} should be extended to contaminated post-consumer cullet with a defined contamination tolerance, and the powder route's finding that irregular milled powder is processable \citep{slmirregular2023} should be tested for laser-filament feedstock drawn from cullet.
  \item \textbf{Fused-silica part scale-up.} The LGD and DGLD groups have shown fully dense walls of tens of millimetres \citep{lgd2024homog,kth2022bonding}; multi-fibre or rod feeding with side-injected heating \citep{capps2021interface} and coaxial ring beams \citep{lgd2021coaxial} are the identified levers to raise deposition rate by an order of magnitude without losing the \SI{20}{nm} surface.
\end{enumerate}

\subsection{Medium-term milestones (3--5 years)}
\begin{enumerate}[leftmargin=*,itemsep=2pt]
  \item \textbf{Borosilicate and aluminosilicate melt extrusion.} Requires a \SIrange{1300}{1450}{\celsius} melter and a nozzle that survives it. The candidates are platinum--rhodium (the Corning application and the fibre-bushing precedent \citep{patent_corning2021,pgm2010}) and dense zirconia; the barrier is cost and rhodium volatilisation. A quick-change platinum-alloy nozzle cartridge of tens of grams, with an inert-gas shroud, is a realistic engineering target and would unlock chemically resistant labware and low-expansion architectural glass. Evenline's stated \SI{1700}{\celsius} G3DP4 variant is the first public commitment \citep{stern2024cgc}.
  \item \textbf{Flow start--stop and multi-nozzle heads.} The MIT patents anticipate shears and multi-orifice nozzles \citep{patent_us9919510}; implementing them lifts the single-continuous-path constraint that currently dictates part design \citep{mit2024masonry} and enables dual-crucible multi-material printing, which the 2026 multi-material review identifies as the frontier \citep{multimat2026}.
  \item \textbf{Hybrid architecture: molten extrusion with laser surface conditioning.} A CO$_2$ laser tracking the nozzle could fire-polish the outer seam while it is still above the softening point, reducing the layer relief that limits optical use, and could locally reheat the previous layer for bonding on cold-chamber machines. Both are routine in polymer and metal directed energy deposition and have no glass-specific obstacle beyond chamber optics.
  \item \textbf{Validated thermo-mechanical simulation.} Coupling a layer-by-layer thermal model (the lumped and conduction models of \cref{sec:physics} are sufficient for soda-lime cords; laser melt pools need the Marangoni-resolved models now appearing \citep{machines2026meltpool,lgd2021coaxial}) to structural relaxation would let the annealing schedule be computed rather than inherited from studio practice, and would give the first predicted residual-stress fields to compare with photoelastic tomography.
  \item \textbf{Structural certification of printed masonry.} A test programme of the scale used for cast-glass masonry at TU Delft, applied to fully hollow printed units, with fire, impact and weathering, would move the 2024 feasibility result \citep{mit2024masonry} into a code-compliant product. This is the shortest path to a building-scale reference project.
\end{enumerate}

\subsection{Longer-term directions}
\emph{Large-format robotic deposition.} Maple's robotic-arm development and Evenline's wide-chamber concept both point to gantry- or arm-mounted heads for fa\c{c}ade panels; the physics favours printing onto preheated float sheet, which supplies the annealing mass and the optical face \citep{stern2024cgc,maple_ammedia}. The chamber then becomes a lehr that travels with the part.

\emph{Programmable optics from instabilities.} The coiling regime of the fluid-mechanical sewing machine \citep{brun2017} and the caustic design methods already applied to printed lamps suggest glass elements whose light distribution is programmed through deposition parameters rather than post-machining---a market that no other material can enter.

\emph{Direct fibre and preform manufacture.} Printed photonic-crystal-fibre preforms drawn to single-mode fibre \citep{laval2023preform} and free-standing fibre printed from SMF-28 \citep{hostetler2018sff} indicate a route from printer to fibre tower that bypasses stack-and-draw; specialty-fibre and chalcogenide IR optics \citep{baudet2019} are the first products.

\emph{Off-world construction.} Molten regolith extrusion from platinum--rhodium crucibles under vacuum \citep{regolithreview2023,regolithvac2024} is a molten-glass printer by another name; the nozzle-corrosion, devitrification and thermal-stress problems of this review are exactly its problems, with the added constraint of solar or laser heating.

\emph{Energy and circularity as the selling proposition.} Remelting cullet consumes less energy than melting virgin batch, printing eliminates moulds, and adhesive-free interlocking units are remeltable at end of life \citep{stern2024cgc,mit2024masonry}. If the crossover batch size (Problem 8) can be shown to be in the tens to hundreds of units for fa\c{c}ade elements, molten glass AM has a business case independent of art and lighting.

\subsection{Recommendations for research funders and practitioners}
\begin{itemize}[leftmargin=*,itemsep=1pt]
  \item Fund the instrumentation gap, not another demonstrator: quantitative stress, transmission and strength characterisation of printed glass is cheaper than a new machine and is what every downstream application needs.
  \item Treat nozzle and crucible corrosion as the critical-path materials problem for melt extrusion; a platinum-alloy or zirconia nozzle life test in three glasses is a well-posed, one-year project.
  \item Connect the two communities: the laser groups have control theory and optics, the extrusion groups have scale and materials handling; a laser-conditioned extrusion head is the natural joint project.
  \item Publish process windows in dimensionless form (Ohnesorge, visco-capillary time, energy density, \(\Delta T\) across the transformation range) so that results transfer between glasses and machines.
  \item Establish a shared benchmark part (a thin wall, a lens blank, a masonry-scale hollow) and report the coupon set above for every new process paper.
\end{itemize}

\subsection{Closing statement}
Molten hot glass AM has reached the point at which the remaining obstacles are engineering and metrology rather than feasibility. The physics is understood well enough to be written down in a page of equations; the machines exist in two complementary forms; the first products are sold. What separates the field from industrial adoption is the absence of quantitative process--property data, the corrosion-limited temperature ceiling of contact extrusion, and the lack of real-time geometric control. Each of these is tractable with methods already proven in adjacent fields, and the roadmap above is a plan to apply them. The complete corpus, figures and source manifest of this review are provided with the LaTeX source so that the next iteration can be built on the same evidence base.
